\section*{Chapter \ref{ch:mx:build-a-matching-engine-and-exchange-simulator} --- Build: A Matching Engine and Exchange Simulator}

\begin{solution}{exo:mx:build-a-matching-engine-and-exchange-simulator:1}
5,013: the packet carries messages 5,001 to 5,012, and heartbeats carry the next expected number without advancing it.
\end{solution}

\begin{solution}{exo:mx:build-a-matching-engine-and-exchange-simulator:2}
Line A 5.3\% ($308/5\,793$), line B 8.9\% ($517/5\,793$), both 0.52\% ($30/5\,793$).
\end{solution}

\begin{solution}{exo:mx:build-a-matching-engine-and-exchange-simulator:3}
The order had already been filled (at 5 seconds), so there was nothing to cancel: reason ``L'', too late to cancel. A cancel of an unknown order gets the
same reason.
\end{solution}

\begin{solution}{exo:mx:build-a-matching-engine-and-exchange-simulator:4}
With one generator in event order, adding a participant or changing one's timing shifts every later draw, so every other participant's latencies and the
feed's losses change: comparisons between runs stop being like for like. Keyed counter-based streams give each purpose its own sequence.
\end{solution}

\begin{solution}{exo:mx:build-a-matching-engine-and-exchange-simulator:5}
To know which incremental messages the snapshot already contains: apply those after its sequence number and skip the ones before. Without it a receiver
applies some messages twice or misses some, and the rebuilt book is wrong in ways that may not show at the top.
\end{solution}

\begin{solution}{exo:mx:build-a-matching-engine-and-exchange-simulator:6}
That the three engines agree on every message, report and snapshot for the fixture's inputs, which exercise every message type and control. It does not
prove that the rules are right (all three could share a misreading of a rule), nor that they agree on inputs the fixture does not contain.
\end{solution}

\begin{solution}{exo:mx:build-a-matching-engine-and-exchange-simulator:7}
Line A loses nothing, line B loses the 166 messages of its outage, and no \omterm{def:mx:build-a-matching-engine-and-exchange-simulator:feed}{retransmission request} is needed: line A delivers every message.
\end{solution}

\begin{solution}{exo:mx:build-a-matching-engine-and-exchange-simulator:8}
The simulator's time goes into agents, the event loop, the background and the feed's packets, not only the engine; and the C++ number includes the
program's start-up and checks, so it is a lower bound on the engine alone. A session's speed-up depends on how much of it is the engine.
\end{solution}

\begin{solution}{pb:mx:build-a-matching-engine-and-exchange-simulator:1}
\textbf{1.} See the definitions.
\textbf{2.} Feed 14, order entry 5, \omterm{def:mx:build-a-matching-engine-and-exchange-simulator:oe}{execution reports} 6, controls 10 (plus SoupBinTCP's own session messages).
\textbf{3.} Accepted; a self-cross trade; an IOC's 200 executed and 300 cancelled; a post-only order resting; an execution at the resting price; an
off-grid rejection; a too-late cancel; an end-of-day expiry.
\textbf{4.} See the definitions.
\textbf{5.} Price-time: in time order at the level; pro rata: Book~1's configurable allocation with a top-order share.
\textbf{6.} The order's limit, the price band, minimum quantities, \omterm{def:mx:build-a-matching-engine-and-exchange-simulator:stp}{self-trade prevention}.
\textbf{7.} Pre-open and opening call, continuous trading or batch auctions, closing call, halted, paused (reopening call); leaving a call uncrosses it.
\textbf{8.} Through the controls R (band) and P (phase), called at scheduled times.
\textbf{9.} Session, first sequence number, count, then messages; A and B carry the same packets.
\textbf{10.} From counter-based streams keyed by line: independent loss, Gilbert--Elliott bursts, duplicates, jitter, outages.
\textbf{11.} Line A lost 308 of 5,793 messages, line B 517; 30 were missing on both; 30 \omterm{def:mx:build-a-matching-engine-and-exchange-simulator:feed}{retransmission requests} filled them.
\textbf{12.} A mid-day snapshot's orders, then the incremental messages after its sequence number; compared with the engine's book ten levels deep.
\textbf{13.} Process both feeds, take each sequence number from whichever arrives first, and recover gaps on both.
\textbf{14.} A simulated clock, fixed tie-breaking, counter-based randomness keyed by purpose, no wall clock in the engine.
\textbf{15.} Four minutes of tape flow plus a chaos agent sending every message type and control: 17,059 journal records.
\textbf{16.} \emph{Named result}: the Python engine passes the thirteen engine-rule scenarios and reproduces the fixture's SHA-256; the C++20 and Rust engines
reproduce the Python engine's output on all 17,059 fixture records byte for byte (and round-trip 4,866 golden messages); on this laptop the Python engine
processes about 60,000 records a second and the C++20 test program more than 500,000.
\textbf{17.} It includes the program's start-up, decoding and byte comparisons.
\textbf{18.} The agents, the latency models, the feed impairments and the live servers' timing.
\textbf{19.} Implied spreads, a FIX gateway, venue-specific rules checked against each venue's documents.
\textbf{20.} So that a difference between two runs can only come from a difference in their inputs.
\end{solution}

\begin{solution}{iq:mx:build-a-matching-engine-and-exchange-simulator:1}
Arbitrate the A and B lines by sequence number; on a gap on both, buffer later messages and request the missing ones from the retransmission service; if
the gap is too old, rebuild from a snapshot and apply the messages after its sequence number.
\iqlookfor{Arbitration, retransmission, snapshots.}
\end{solution}

\begin{solution}{iq:mx:build-a-matching-engine-and-exchange-simulator:2}
A map from price to level (a sorted structure or an array indexed by ticks), each level a FIFO of orders, and a hash from order reference to its node
so a cancel is O(1); keep the best prices cached.
\iqlookfor{Levels, FIFO, reference index.}
\end{solution}

\begin{solution}{iq:mx:build-a-matching-engine-and-exchange-simulator:3}
A discrete-event loop on a simulated clock with deterministic tie-breaking, counter-based random streams keyed by purpose, no wall-clock or hash-order
dependence, and a journal of every input; test by hashing outputs of repeated runs.
\iqlookfor{Clock, randomness, ordering, testing.}
\end{solution}

\begin{solution}{iq:mx:build-a-matching-engine-and-exchange-simulator:4}
An acceptance, one execution per fill (with the quantity left), then a cancel for the remainder with the IOC reason; the feed shows the executions against
the resting orders.
\iqlookfor{Report sequence; leaves quantity.}
\end{solution}

\begin{solution}{iq:mx:build-a-matching-engine-and-exchange-simulator:5}
\omterm{def:mx:build-a-matching-engine-and-exchange-simulator:stp}{Cancel on disconnect} prevents stale orders from trading while a participant cannot manage them; \omterm{def:mx:build-a-matching-engine-and-exchange-simulator:stp}{self-trade prevention} prevents wash trades between a firm's
own strategies and the fees and reporting problems they cause.
\iqlookfor{Stale orders; wash trades.}
\end{solution}

\begin{solution}{iq:mx:build-a-matching-engine-and-exchange-simulator:6}
Replay the journal to the first record whose output differs, decode both outputs, read the rule in the specification, and write a minimal scenario that
reproduces it; add it to the fixture so the three engines are held to it.
\iqlookfor{Bisection on the journal; a new test.}
\end{solution}
